\chapter{Tests}

The test suite uses GoogleTest and is configured in the \texttt{test/} folder. The goal is to ensure each module reacts correctly to nominal cases and edge cases, with particular attention to safety conditions.

\section{Structure}
The \texttt{CMakeLists.txt} file in \texttt{test/} downloads GoogleTest via FetchContent.
The generated executable is called \texttt{pinacoteca\_tests} and Arduino API mocks are in \texttt{test/mocks/}.
The test entry point is \texttt{all\_tests.cpp}, where each feature group has dedicated tests.

\section{Execution}
Two modes are available:
\begin{itemize}
  \item Direct CMake usage:
  \begin{lstlisting}[style=cppstyle]
cd test
cmake -S . -B build
cmake --build build -j
ctest --test-dir build --output-on-failure
  \end{lstlisting}
  \item \texttt{run\_tests.sh} script that manages build, clean, and Arduino compilation.
\end{itemize}

\section{Test rationale}
Below is the detailed explanation of the checks, with objectives, conditions, and expectations.

\subsection{LED}
\begin{itemize}
  \item \textbf{Objective}: ensure output primitives write to the correct pins.
  \item \textbf{Conditions}: direct calls to \texttt{led} and \texttt{ledDimming} with known values.
  \item \textbf{Expectations}: \texttt{digitalWrite} and \texttt{analogWrite} receive the requested value.
  \item \textbf{Tests}:
  \begin{tabular}[t]{@{}l@{}}
    \texttt{LedWriteDigitalState}\\
    \texttt{LedDimmingWritesPwm}
  \end{tabular}
\end{itemize}

\subsection{Stoplight}
\begin{itemize}
  \item \textbf{Objective}: validate green/red semantics based on capacity.
  \item \textbf{Conditions}: counts below and equal to the maximum.
  \item \textbf{Expectations}: green on below capacity, red on at full capacity.
  \item \textbf{Tests}:
  \begin{tabular}[t]{@{}l@{}}
    \texttt{StoplightBelowMaxSetsGreen}\\
    \texttt{StoplightAtMaxSetsRed}
  \end{tabular}
\end{itemize}

\subsection{Servomotor}
\begin{itemize}
  \item \textbf{Objective}: prevent movement beyond mechanical limits.
  \item \textbf{Conditions}: commands with valid angles and angles outside 0--180.
  \item \textbf{Expectations}: movement allowed within range, blocked and logged outside range.
  \item \textbf{Tests}:
  \begin{tabular}[t]{@{}l@{}}
    \texttt{AntiSufferingServoAllowsSafeMove}\\
    \texttt{AntiSufferingServoBlocksOutOfRange}
  \end{tabular}
\end{itemize}

\subsection{Turnstile}
\begin{itemize}
  \item \textbf{Objective}: count entrances/exits correctly and respect the maximum.
  \item \textbf{Conditions}: simulated IN/OUT button states; variable max people.
  \item \textbf{Expectations}: coherent increment and decrement, no limit overruns.
  \item \textbf{Tests}:
  \begin{tabular}[t]{@{}l@{}}
    \texttt{TurnstileIncrementsAndDecrementsPeople}\\
    \texttt{TurnstileRespectsMaxPeople}\\
    \texttt{TurnstileBeginConfiguresPins}
  \end{tabular}
\end{itemize}

\subsection{Temperature}
\begin{itemize}
  \item \textbf{Objective}: handle abnormal NTC sensor readings.
  \item \textbf{Conditions}: analog values at extremes (0 and 1023) and an intermediate value.
  \item \textbf{Expectations}: sentinel \texttt{-999.0} at extremes, finite value under normal conditions.
  \item \textbf{Tests}:
  \begin{tabular}[t]{@{}l@{}}
    \texttt{TemperatureReaderReturnsErrorAtExtremes}\\
    \texttt{TemperatureReaderReturnsFiniteForValidInput}
  \end{tabular}
\end{itemize}

\subsection{Thermostat}
\begin{itemize}
  \item \textbf{Objective}: verify heating/cooling logic and fail-safe behavior.
  \item \textbf{Conditions}: valid reading above/below target, and invalid reading.
  \item \textbf{Expectations}: correct pin activation, outputs off on sensor error.
  \item \textbf{Tests}:
  \begin{tabular}[t]{@{}l@{}}
    \texttt{ThermostatHeatingBranch}\\
    \texttt{ThermostatCoolingBranch}\\
    \texttt{ThermostatSensorErrorTurnsEverythingOff}\\
    \texttt{ThermostatBeginSetsPinModesAndSetGetTarget}
  \end{tabular}
\end{itemize}

\subsection{Humidity and HumidifierControl}
\begin{itemize}
  \item \textbf{Objective}: validate DHT22 reading and threshold logic.
  \item \textbf{Conditions}: high humidity, low humidity, and invalid reading (\texttt{NAN}).
  \item \textbf{Expectations}: actuator consistent with threshold and off on error.
  \item \textbf{Tests}:
  \begin{tabular}[t]{@{}l@{}}
    \texttt{HumidityReadAndError}\\
    \texttt{HumidifierHighHumidityTurnsOnPin}\\
    \texttt{HumidifierLowHumidityTurnsOffPin}\\
    \texttt{HumidifierSensorErrorReturnsFalse}\\
    \texttt{HumidifierSetGetTargetWorks}
  \end{tabular}
\end{itemize}

\subsection{Photoresistor and LightingControl}
\begin{itemize}
  \item \textbf{Objective}: ensure conversion and PWM align with the target.
  \item \textbf{Conditions}: analog values at extremes and intermediate values.
  \item \textbf{Expectations}: lux clamped at the edges and coherent PWM (255 with low light, 0 with high light).
  \item \textbf{Tests}:
  \begin{tabular}[t]{@{}l@{}}
    \texttt{PhotoresistorLuxEdgeCases}\\
    \texttt{LightingControlWritesPwm}\\
    \texttt{LightingControlBeginAndSetGetTarget}
  \end{tabular}
\end{itemize}

\subsection{DisplayInterface}
\begin{itemize}
  \item \textbf{Objective}: verify the abstract LCD layer.
  \item \textbf{Conditions}: calls to \texttt{begin}, \texttt{clear}, \texttt{setCursor}, \texttt{print}.
  \item \textbf{Expectations}: the mock LCD buffer contains the expected strings.
  \item \textbf{Tests}:
  \begin{tabular}[t]{@{}l@{}}
    \texttt{DisplayInterfaceBasicOperations}
  \end{tabular}
\end{itemize}

\subsection{DisplayPanel}
\begin{itemize}
  \item \textbf{Objective}: verify page rotation and diagnostic messages.
  \item \textbf{Conditions}: simulated time advancement and sensor value changes.
  \item \textbf{Expectations}: date/time page, status page, and error page shown in rotation.
  \item \textbf{Tests}:
  \begin{tabular}[t]{@{}l@{}}
    \texttt{DisplayPanelRotatesPages}\\
    \texttt{DisplayDatePageHasValidWeekdayPrefix}\\
    \texttt{DisplayStatusPageShowsTurnstileCountExactly}\\
    \texttt{DisplayErrorPageShowsHumidityHigh}\\
    \texttt{DisplayErrorPageShowsTurnstileError}
  \end{tabular}
\end{itemize}

\section{run\_tests.sh script}
Main options:
\begin{itemize}
  \item \texttt{--clean}: deletes the build folder.
  \item \texttt{--with-arduino}: compiles the sketch with arduino-cli.
  \item \texttt{--upload}: uploads the sketch to a board (requires \texttt{--port}).
  \item \texttt{--board <fqbn>}: selects the board, default \texttt{arduino:avr:uno}.
\end{itemize}
